\section{Введение}

\subsection*{Машинное и глубокое обучение как инструмент порождающего моделирования}

Машинное обучение~--- область прикладной математики, посвящённая восстановлению зависимостей по эмпирическим данным~\cite{vapnik1999overview, hastie2009elements}. Формально пусть~$\mathcal{X}\subset\mathbb{R}^d$~--- пространство объектов, $\mathcal{Y}$~--- пространство ответов, $\pi$~--- неизвестное совместное распределение на~$\mathcal{X}\times\mathcal{Y}$. Параметрическая модель~$f_\theta:\mathcal{X}\to\mathcal{Y}$, $\theta\in\Theta$, обучается минимизацией эмпирического риска
\[
    \hat{\theta} = \argmin_{\theta\in\Theta}\;\frac{1}{n}\sum_{i=1}^n \ell\bigl(f_\theta(\mathbf{x}_i), y_i\bigr), \qquad (\mathbf{x}_i,y_i)\overset{\text{iid}}{\sim}\pi,
\]
аппроксимирующего теоретический риск~$R(\theta)=\mathbb{E}_{(\mathbf{x},y)\sim\pi}\,\ell(f_\theta(\mathbf{x}),y)$. Глубокое обучение~\cite{lecun2015deep, goodfellow2016deep} реализует~$f_\theta$ как многослойную нейронную сеть~--- композицию параметризованных аффинных преобразований и нелинейностей; теоремы об универсальной аппроксимации~\cite{cybenko1989approximation, hornik1989multilayer, leshno1993multilayer} гарантируют, что класс таких функций плотен в~$C(\mathcal{X})$ относительно равномерной нормы при минимальных условиях на нелинейность. Сочетание стохастического градиентного спуска~\cite{robbins1951stochastic, kingma2014adam} с алгоритмом обратного распространения ошибки~\cite{rumelhart1986learning} и параллельных вычислений на графических ускорителях позволило применять модели с десятками миллионов параметров к задачам компьютерного зрения~\cite{krizhevsky2012imagenet, he2016deep}, обработки естественного языка~\cite{vaswani2017attention, devlin2018bert} и порождающего моделирования.

В последнем направлении центральной задачей является аппроксимация распределения~$\pi_\mathrm{data}$ либо его сэмплера. Классическая постановка состоит в построении параметрической плотности~$p_\theta$, минимизирующей расхождение Кульбака--Лейблера $\mathrm{KL}(\pi_\mathrm{data}\,\|\,p_\theta)$~\cite{kingma2013auto, rezende2014stochastic}. Альтернативный неявный подход~--- обучение порождающего отображения~$G_\theta:\mathcal{Z}\to\mathcal{X}$, при котором $G_{\theta\#}\pi_\mathrm{prior}\approx\pi_\mathrm{data}$, как в генеративно-состязательных сетях~\cite{goodfellow2014generative, arjovsky2017wasserstein}. Современный сдвиг парадигмы связан с переходом от одношагового порождения к итеративному, основанному на стохастических дифференциальных уравнениях~--- диффузионных моделях~\cite{sohl2015deep, ho2020denoising, song2020score} и моделях согласования оценок (score matching~\cite{hyvarinen2005estimation, vincent2011connection}). Их теоретическое обоснование лежит в области стохастического анализа, теории Фоккера--Планка и оптимального транспорта.

Качественный скачок 2022--2023~годов~--- появление флоу матчинга~\cite{lipman2022flow, albergo2022building, liu2022flow}, который заменяет стохастический обратный процесс детерминированным потоком~$\phi_t:\mathbb{R}^d\to\mathbb{R}^d$, заданным обыкновенным дифференциальным уравнением
\[
    \frac{\mathrm{d}\phi_t(\mathbf{x})}{\mathrm{d}t} = u_t(\phi_t(\mathbf{x})), \qquad \phi_0=\mathrm{id},
\]
с маргинальным полем скоростей~$u_t$, согласованным с уравнением непрерывности $\partial_t\pi_t+\nabla\!\cdot\!(\pi_t u_t)=0$~\cite{villani2008optimal}. Поле~$u_t$ аппроксимируется нейронной сетью~$f_\theta$ через симуляционно-свободный функционал, эквивалентный по градиентам $\ell_2$-регрессии на условную скорость~\cite{lipman2022flow}. Связь с минимизацией кинетической энергии Бенаму--Бренье~\cite{benamou2000computational} и теоремой Бренье~\cite{brenier1991polar} обеспечивает геометрическую интерпретацию: при линейном интерполянте траектории спрямляются, а сложность задачи аппроксимации падает. Практическое следствие~--- инференс за один шаг численного интегрирования вместо $T\sim10^3$ шагов диффузии.

\subsection*{Прикладная задача: трансляция КТ-контраста}

Изложенный математический аппарат естественно применяется к одной из наиболее остро стоящих задач медицинской визуализации~--- порождению парных контрастных и нативных компьютерных томограмм. Компьютерная томография занимает центральное место в современной клинической визуализации брюшной полости: ежегодно в США выполняется свыше 93~миллионов КТ-исследований~\cite{10.1001/jamainternmed.2025.0505}, при этом стандартные диагностические протоколы нередко предусматривают сразу несколько серий~--- нативную (без контрастного усиления) и контрастную артериальную фазы, приобретаемые с интервалом 15--35 секунд после внутривенного введения йодсодержащего вещества~\cite{MayoSmith2017ManagementOI}. Каждая из фаз несёт самостоятельную диагностическую ценность: нативная серия необходима для оценки плотности тканей в единицах Хаунсфилда, выявления кальцификатов и установления базового фона; контрастная~--- для чёткой визуализации сосудистых структур и паренхиматозных органов.

Публично доступные наборы медицинских данных редко содержат обе фазы. Крупнейший на сегодняшний день датасет AbdomenAtlas~\cite{li2024abdomenatlas} (8\,448 томограмм) преимущественно ограничивается контрастной фазой. Отсутствие нативной серии делает невозможным ряд диагностических задач~--- от подсчёта кальций-скоринга до оценки плотностных характеристик тканей. Получение дополнительной КТ-серии влечёт лучевую нагрузку, ограничивает пропускную способность оборудования и сопряжено с риском нефропатии, индуцированной контрастным веществом~\cite{doi:10.1148/radiol.2392050413}.

Совокупность этих обстоятельств формирует практический запрос: построить алгоритм трансляции между нативными и контрастными КТ. Решение задачи позволит расширять обучающие наборы данных без проведения дополнительных исследований. Особенностью задачи, отличающей её от большинства приложений синтеза медицинских изображений, является требование точного воспроизведения абсолютных значений плотности в HU~--- стандартизованной радиологической шкалы, где диагностические пороги (например, 70~HU для кальцификации) требуют численного, а не лишь визуального, соответствия.

Применение методов глубокого обучения к данной задаче исследовалось в нескольких направлениях: регрессионные сети~\cite{article_dual}, генеративно-состязательные сети~\cite{hu2021bidirectional, article_nature}, диффузионные модели~\cite{article_nature, kim2024adaptive} и специализированный метод SMILE~\cite{liu2025see}. Диффузионные подходы демонстрируют высокое качество, однако страдают от вычислительно затратного многошагового вывода. Применение флоу матчинга~\cite{lipman2022flow, albergo2022building} к медицинской визуализации остаётся малоизученным~\cite{yazdani2025flow, reynaud2025echoflow}; ни одна из известных работ не рассматривала его для задачи парной трансляции КТ-контраста, при которой пары $(\mathbf{x}_N, \mathbf{x}_C)$ от одного пациента дают эмпирическое приближение оптимального транспортного сочленения.

\paragraph{Постановка задачи и вклад работы.}
В данной работе исследуется математическая задача построения непрерывного потока, согласующего распределения нативных~$\pi_N$ и контрастных~$\pi_C$ КТ-срезов, в рамках теории условного флоу матчинга. Предлагается метод \textbf{Medical Flow Matching~(MFM)} и сопутствующая нейросетевая параметризация \textbf{TimeResNet}, служащая инструментом экспериментальной проверки теоретических утверждений о (а)~симуляционно-свободном обучении, (б)~оптимальности одношагового интегрирования при линейном интерполянте и (в)~двунаправленности трансляции, вытекающей из симметрии интерполянта.

Основные результаты работы:
\begin{enumerate}
    \item Разработан конвейер условного флоу матчинга для двунаправленной трансляции КТ-контраста, требующий лишь однократного вычисления сети при инференсе и обеспечивающий ускорение в \textbf{354~раза} относительно DDPM (1000 шагов).
    \item Предложена архитектура TimeResNet с (а)~пошаговым синусоидальным внедрением временного сигнала через аффинную модуляцию в каждый блок и (б)~свёрточным модулем самовнимания на боттлнеке; абляционное исследование подтверждает вклад каждого компонента (+1.71 и +1.77\,HU MAE при удалении).
    \item Экспериментально подтверждено, что единственная сеть~$f_\theta$ работает в обоих направлениях трансляции, что вдвое сокращает затраты на обучение. Сгенерированные изображения пригодны для обучения последующих сегментационных моделей: Dice~=~0.926 против 0.966 на реальных данных.
\end{enumerate}

\paragraph{Структура работы.}
В разделе~\ref{sec:related} приводится обзор релевантных публикаций, охватывающий методы машинного обучения для синтеза медицинских изображений, задачу трансляции КТ-контраста и методологию флоу матчинга. Раздел~\ref{sec:notation} фиксирует обозначения, формулирует математическую постановку задачи и излагает теоретические основы: уравнение непрерывности, связь с оптимальным транспортом, теорему симуляционно-свободного обучения, сравнение с диффузионными моделями и численные методы интегрирования ОДУ. Раздел~\ref{sec:method} содержит описание метода MFM и архитектуры TimeResNet. В разделе~\ref{sec:experiments} представлены данные, детали обучения, метрики качества, количественные и качественные результаты, абляционные эксперименты и анализ применимости для задачи сегментации. Раздел~\ref{sec:conclusion} подводит итог проделанной работы и намечает направления дальнейших исследований.
